\documentclass[12pt]{article}
\usepackage[T1]{fontenc}
\usepackage[utf8]{inputenc}
\usepackage[polish]{babel}
\usepackage{lmodern}
\usepackage[a4paper,margin=2.4cm]{geometry}
\usepackage{xcolor}
\usepackage{graphicx}
\usepackage{amsmath,amssymb}
\usepackage{booktabs}
\usepackage{array}
\usepackage{enumitem}
\usepackage{hyperref}
\usepackage{tcolorbox}
\usepackage{fancyhdr}
\definecolor{WEIIPrimary}{RGB}{0,130,150}
\definecolor{WEIIPrimaryDark}{RGB}{0,86,100}
\definecolor{WEIIPrimaryLight}{RGB}{229,246,248}
\definecolor{PRzNavy}{RGB}{31,63,142}
\definecolor{RLTextDark}{RGB}{34,34,34}
\definecolor{RLGray}{RGB}{90,90,90}
\definecolor{RLLightGray}{RGB}{245,246,248}
\definecolor{RLWarning}{RGB}{184,86,0}
\definecolor{RLExample}{RGB}{52,112,64}
\definecolor{RLCodeBg}{RGB}{248,248,248}

\usepackage{listings}
\usepackage{xcolor}

\lstdefinestyle{rlpython}{%
    language=Python,
    basicstyle=\ttfamily\scriptsize,
    keywordstyle=\color{WEIIPrimaryDark}\bfseries,
    commentstyle=\color{RLGray}\itshape,
    stringstyle=\color{RLExample},
    numbers=left,
    numberstyle=\tiny\color{RLGray},
    numbersep=6pt,
    showstringspaces=false,
    breaklines=true,
    tabsize=4,
    frame=single,
    rulecolor=\color{WEIIPrimary!35},
    backgroundcolor=\color{RLCodeBg},
    captionpos=b
}

\lstset{style=rlpython}


\hypersetup{colorlinks=true,linkcolor=WEIIPrimaryDark,urlcolor=WEIIPrimaryDark,citecolor=WEIIPrimaryDark}
\setlength{\headheight}{15pt}
\pagestyle{fancy}
\fancyhf{}
\lhead{Uczenie się ze wzmocnieniem}
\rhead{WEiI PRz}
\cfoot{\thepage}

\newtcolorbox{taskbox}[1]{colback=WEIIPrimaryLight,colframe=WEIIPrimary,title=\textbf{#1}}
\newtcolorbox{notebox}[1]{colback=RLLightGray,colframe=RLGray,title=\textbf{#1}}
\newtcolorbox{warninghandout}[1]{colback=RLWarning!8,colframe=RLWarning,title=\textbf{#1}}

\newcommand{\handouttitle}[3]{%
\begin{center}
{\Large\bfseries #1}\\[0.2cm]
{\large #2}\\[0.15cm]
{\small #3}
\end{center}
\vspace{0.4cm}
\hrule
\vspace{0.6cm}
}

\newcommand{\EnableLabNumbering}{%
  \renewcommand{\thesection}{L\arabic{section}}%
  \renewcommand{\thesubsection}{\thesection.\arabic{subsection}}%
  \renewcommand{\thesubsubsection}{\thesubsection.\arabic{subsubsection}}%
}


\title{Raport projektowy\\Uczenie się ze wzmocnieniem}
\author{Patryk Kubik}
\date{2026/2027}

\begin{document}
\maketitle
\tableofcontents
\newpage

\section{Informacje podstawowe}

\begin{center}
\begin{tabular}{p{0.34\textwidth}p{0.56\textwidth}}
\toprule
Element & Informacja \\
\midrule
Imię i nazwisko & Patryk Kubik \\
Numer albumu & 179538 \\
Kierunek/specjalność & Informatyka/Sztuczna Inteligencja \\
Tytuł projektu & Zastosowanie uczenia ze wzmocnieniem w grze typu Flappy Bird \\
Link do repozytorium & \url{https://github.com/PATER10/flappy-bird-reinforcement-learning-project} \\
Data oddania & Do uzupełnienia \\
\bottomrule
\end{tabular}
\end{center}

\section{Temat i cel projektu}

\subsection{Opis problemu}

Projekt polega na rozwiązaniu problemu nieskończonego horyzontu (survivalowego) w dynamicznym środowisku przewijanym bocznie (2D). Agent sterujący obiektem (ptakiem) musi poruszać się w przestrzeni, podlegając wirtualnej grawitacji, i unikać kolizji z pojawiającymi się na jego drodze przeszkodami (rurami). Zderzenie z dowolną przeszkodą lub podłożem kończy próbę.

\subsection{Cel agenta}

Celem agenta jest maksymalizacja skumulowanej nagrody, co przekłada się na przetrwanie jak najdłużej bez kolizji z przeszkodami i zdobycie jak najwyższego wyniku punktowego. Działanie agenta uznaje się za skuteczne, gdy potrafi on w pełni autonomicznie, bez interwencji człowieka, omijać przeszkody i utrzymywać się w wyznaczonych granicach środowiska przez określoną z góry liczbę wygenerowanych rur, zapobiegając szybkiemu zakończeniu epizodu.

\subsection{Uzasadnienie zastosowania uczenia się ze wzmocnieniem}

Problem gry Flappy Bird stanowi klasyczny przykład procesu decyzyjnego, w którym uczenie się ze wzmocnieniem (RL) ma ogromne zastosowanie, ponieważ:
\begin{itemize}
    \item agent musi sekwencyjnie i w czasie rzeczywistym podejmować decyzję o wykonaniu skoku lub powstrzymaniu się od niego na podstawie zmieniających się współrzędnych przeszkód,
    \item każda decyzja (np. skok) bezpośrednio wpływa na przyszły stan środowiska -- zbyt wczesny lub zbyt późny skok może doprowadzić do nieuchronnej kolizji w kolejnych klatkach symulacji,
    \item istotny jest wynik długoterminowy, a nie natychmiastowa nagroda -- pojedynczy skok może nie przynieść od razu punktu (nagrody), ale jest niezbędny do uniknięcia uderzenia w rurę oddaloną o kilkadziesiąt pikseli,
    \item problem ten cechuje się bardzo dużą, ciągłą przestrzenią stanów. Zastosowanie agenta uczącego się (w szczególności metod Deep RL) jest uzasadnione, ponieważ agent jest w stanie sam wywnioskować optymalną politykę sterowania, bez konieczności ręcznego programowania reguł ,,jeżeli ptak jest $x$ centymetrów od rury, to skacz'', które byłyby nieskuteczne.
\end{itemize}

\section{Formalizacja problemu RL/MDP}

\subsection{Agent i środowisko}

Agentem w tym problemie jest model algorytmiczny (oparty na głębokiej sieci neuronowej), który odpowiada za sterowanie ptakiem. Środowiskiem jest fizyczny silnik gry stworzony w Unity, który zarządza grawitacją, ruchem obiektu i generowaniem przeszkód. W jednym kroku epizodu:
\begin{enumerate}
    \item Środowisko mierzy aktualne współrzędne ptaka i rur, po czym przesyła te dane (stan) do agenta poprzez interfejs API.
    \item Agent na podstawie otrzymanych danych liczbowych podejmuje decyzję o akcji.
    \item Agent wysyła akcję z powrotem do środowiska.
    \item Środowisko (Unity) wykonuje ruch, sprawdza kolizje i wylicza nagrodę, po czym cały proces się powtarza dla nowej klatki animacji.
\end{enumerate}

\subsection{Stan lub obserwacja}

W tym przypadku agent otrzymuje obserwację, która bardzo dokładnie odwzorowuje rzeczywisty stan środowiska ważny do podjęcia decyzji.

\begin{center}
\begin{tabular}{p{0.30\textwidth}p{0.60\textwidth}}
\toprule
Składowa stanu/obserwacji & Znaczenie \\
\midrule
\texttt{bird\_y} & Współrzędna Y położenia ptaka na ekranie. \\
\texttt{bird\_velocity} & Aktualna prędkość opadania lub wznoszenia ptaka. \\
\texttt{pipe\_dist\_x} & Odległość w poziomie do najbliższej nadciągającej rury. \\
\texttt{pipe\_gap\_y} & Współrzędna Y środka wolnej przestrzeni (szczeliny) pomiędzy dolną a górną rurą. \\
\bottomrule
\end{tabular}
\end{center}

\subsection{Akcje}

Zbiór akcji dostępnych dla agenta jest dyskretny i bardzo mały. W każdym kroku decyzyjnym agent ma do dyspozycji jedną z dwóch akcji. Żadna z nich nie jest zablokowana w zależności od stanu.

\begin{center}
\begin{tabular}{p{0.30\textwidth}p{0.60\textwidth}}
\toprule
Akcja & Opis \\
\midrule
\texttt{0} (Brak akcji) & Ptak nie podejmuje żadnego działania, opadając swobodnie pod wpływem grawitacji. \\
\texttt{1} (Skok) & Agent uruchamia impuls siły działający na ptaka pionowo w górę, nadając mu natychmiastową prędkość wznoszenia. \\
\bottomrule
\end{tabular}
\end{center}

\subsection{Funkcja nagrody}

Funkcja nagrody została zaprojektowana w taki sposób, aby promować przetrwanie i unikanie przeszkód. 

\begin{equation}
    r_t = \begin{cases} 
    -100 & \text{za kolizję (rura, ziemia, wyjście poza ekran),} \\
    +10 & \text{za pomyślny przelot przez szczelinę pomiędzy rurami,} \\
    +0.1 & \text{za każdy przeżyty krok (brak kolizji i brak przelotu).}
    \end{cases}
\end{equation}

W tej sekcji wyjaśniono:
\begin{itemize}
    \item \textbf{za co agent otrzymuje nagrodę dodatnią:} Agent otrzymuje dużą nagrodę dodatnią (+10) za bezpieczne pokonanie przeszkody (przekroczenie współrzędnej X rury). Otrzymuje również bardzo małą nagrodę (+0.1) ułamkową za samo utrzymanie się przy życiu w każdej klatce symulacji.
    \item \textbf{za co agent otrzymuje karę:} Agent otrzymuje dotkliwą karę (-100) za zderzenie z rurą, uderzenie w podłoże, a także za wylecenie poza górną krawędź ekranu.
    \item \textbf{czy występuje kara za krok:} W klasycznym ujęciu, gdzie celem jest jak najszybsze dotarcie do mety, stosuje się ujemną karę za krok. Ponieważ w grze Flappy Bird celem jest przetrwanie (nieskończony horyzont), zamiast ujemnej kary za krok zastosowano minimalną nagrodę dodatnią za krok, zachęcającą agenta do unikania stanów terminalnych.
    \item \textbf{czy występują ograniczenia twarde lub miękkie:} Ograniczenia twarde stanowią fizyczne kolidery środowiska -- obiekty rur oraz podłoże. Ich naruszenie jest niedozwolone i natychmiast kończy epizod porażką.
    \item \textbf{czy nagroda nie zachęca agenta do niepożądanego zachowania:} Funkcja nagrody uwzględnia karę za opuszczenie obszaru gry od góry (wylecenie poza ekran). Zapobiega to niepożądanemu zachowaniu, polegającemu na ciągłym skakaniu, aby przelecieć bezpiecznie wysoko nad rurami, omijając w ten sposób przeszkody.
\end{itemize}

\subsection{Warunki zakończenia epizodu}

Epizod (jedna próba przejścia gry) ulega natychmiastowemu zakończeniu po wejściu środowiska w stan terminalny. Następuje to w przypadku spełnienia co najmniej jednego z poniższych warunków:
\begin{itemize}
    \item \textbf{Kolizja z przeszkodą:} Punkt kolizyjny ptaka wejdzie w interakcję z prostokątem kolizyjnym dolnej lub górnej rury.
    \item \textbf{Kolizja z podłożem:} Ptak opadnie na poziom ziemi.
    \item \textbf{Opuszczenie obszaru gry:} Ptak wyleci poza górną krawędź ekranu (zabezpieczenie przed omijaniem rur górą).
    \item \textbf{Osiągnięcie celu (limit rur):} Agent z powodzeniem pokona 100 rur bez kolizji. Jest to sztuczny warunek zakończenia wprowadzony w celu uniknięcia nieskończonych epizodów w przypadku wyuczenia się przez agenta optymalnej polityki sterowania.
\end{itemize}

\subsection{Założenia i ograniczenia środowiska}

W ramach projektu przyjęto następujące założenia i uproszczenia modelu środowiska:
\begin{itemize}
    \item \textbf{Reprezentacja stanu:} Zrezygnowano z analizy obrazu (pikseli) na rzecz bezpośredniego przekazywania dokładnych współrzędnych przestrzennych. Agent ,,nie widzi'' gry, lecz analizuje wektor liczb. Znacząco zmniejsza to złożoność obliczeniową i pozwala na wykorzystanie prostszej architektury sieci.
    \item \textbf{Uproszczona fizyka:} Model fizyczny środowiska uwzględnia jedynie stałą siłę grawitacji (działającą wzdłuż osi Y) oraz skokowe impulsy siły (akcja skoku). Pominęto efekty aerodynamiczne, takie jak opór powietrza, rotację obiektu czy zmienny wiatr.
    \item \textbf{Komunikacja przez API:} Ponieważ agent (Python) i środowisko (Unity) komunikują się przez zewnętrzny interfejs, pomija się kwestię mikrosekundowych opóźnień sieciowych (tzw. lagów). Zakłada się, że przy symulacji uruchomionej w przyspieszonym czasie wymiana pakietów stanu i akcji zachodzi w sposób w pełni synchroniczny z klatkami fizyki.
\end{itemize}

\section{Implementacja}

\subsection{Struktura repozytorium}

Należy krótko opisać strukturę repozytorium i najważniejsze pliki projektu.

\begin{verbatim}
repozytorium/
|-- src/
|-- experiments/
|-- report/
|-- results/
`-- README.md
\end{verbatim}

\subsection{Środowisko}

Należy opisać sposób implementacji środowiska, w szczególności funkcje inicjalizacji epizodu oraz wykonania kroku środowiska.

\subsection{Agent bazowy}

Należy opisać agenta losowego, heurystycznego albo inny baseline wykorzystany do porównania z agentem uczącym się.

\subsection{Agent uczący się}

Należy opisać zastosowany algorytm RL lub DRL, jego parametry oraz sposób implementacji. W przypadku Q-tab należy opisać reprezentację tablicy wartości. W przypadku DRL należy opisać wejścia modelu, wyjścia modelu, architekturę oraz sposób uczenia.

\section{Plan eksperymentów}

Należy opisać eksperymenty zaplanowane w celu oceny działania agenta.

\begin{center}
\begin{tabular}{p{0.24\textwidth}p{0.30\textwidth}p{0.30\textwidth}}
\toprule
Eksperyment & Co jest zmieniane? & Co jest mierzone? \\
\midrule
E1 & Do uzupełnienia & Do uzupełnienia \\
E2 & Do uzupełnienia & Do uzupełnienia \\
E3 & Do uzupełnienia & Do uzupełnienia \\
\bottomrule
\end{tabular}
\end{center}

W tej sekcji należy wskazać:
\begin{itemize}
    \item liczbę epizodów uczenia,
    \item liczbę powtórzeń eksperymentu albo uzasadnienie braku powtórzeń,
    \item parametry uczenia,
    \item metryki oceny,
    \item warianty środowiska, map lub funkcji nagrody,
    \item sposób porównania z agentem bazowym.
\end{itemize}

\section{Wyniki}

W tej sekcji należy przedstawić wyniki eksperymentów. Wyniki powinny być pokazane w postaci tabel, wykresów lub czytelnych zestawień.

\begin{center}
\begin{tabular}{p{0.28\textwidth}p{0.20\textwidth}p{0.20\textwidth}p{0.20\textwidth}}
\toprule
Wariant & Średni zwrot & Skuteczność & Długość epizodu \\
\midrule
Agent bazowy & Do uzupełnienia & Do uzupełnienia & Do uzupełnienia \\
Agent RL/DRL & Do uzupełnienia & Do uzupełnienia & Do uzupełnienia \\
\bottomrule
\end{tabular}
\end{center}

\begin{figure}[h]
    \centering
    \fbox{\parbox{0.82\textwidth}{\centering Miejsce na wykres procesu uczenia}}
    \caption{Przykładowy wykres procesu uczenia agenta.}
\end{figure}

\section{Analiza i wnioski}

Należy zinterpretować wyniki. W tej części należy odpowiedzieć m.in. na pytania:
\begin{itemize}
    \item Czy agent uczący się działa lepiej niż agent bazowy?
    \item Jak zmieniało się zachowanie agenta w trakcie uczenia?
    \item Które elementy środowiska były najtrudniejsze?
    \item Czy funkcja nagrody prowadziła do oczekiwanego zachowania?
    \item Czy wystąpiły przypadki niepożądanego zachowania agenta?
    \item Jakie są najważniejsze ograniczenia rozwiązania?
\end{itemize}

\section{Instrukcja uruchomienia}

Należy opisać sposób uruchomienia projektu od czystego repozytorium. Przykładowo:

\begin{verbatim}
git clone https://github.com/...
cd nazwa-repozytorium
python -m venv .venv
.venv\Scripts\activate
pip install -r requirements.txt
python src/main.py
\end{verbatim}

Należy również wskazać, jak odtworzyć eksperymenty oraz gdzie zapisywane są wyniki.

\section{Wsparcie AI}

W tej sekcji należy obowiązkowo opisać wykorzystanie narzędzi sztucznej inteligencji. Jeżeli student nie korzystał z takich narzędzi, należy wpisać jednoznaczną informację, że nie korzystano ze wsparcia AI.

\begin{center}
\begin{tabular}{p{0.13\textwidth}p{0.19\textwidth}p{0.21\textwidth}p{0.27\textwidth}}
\toprule
Data & Narzędzie/model & Cel użycia & Transkrypt / link / zakres \\
\midrule
26/10/2026 - 05/11/2026 & Gemini 3.1 PRO & Konsultacja w celu wybrania tematu projektu, pomoc w pisaniu 3 pierwszych rozdziałów raportu & https://share.gemini.google/KI2bDKycjaVm / Wszystkie treści zostały dogłębnie przeanalizowane i dostosowane do wymagań, zgodne opisy zostały ujęte w raporcie \\
\bottomrule
\end{tabular}
\end{center}

Należy zapisać transkrypty rozmów z modelami AI wraz z datą oraz nazwą modelu. Jeżeli rozmowa odbywała się poza narzędziami programistycznymi, należy dodatkowo udostępnić całą rozmowę poprzez link, o ile użyte narzędzie umożliwia takie udostępnienie.

Student powinien także wyjaśnić, które elementy wygenerowane lub zasugerowane przez AI zostały samodzielnie sprawdzone, zmodyfikowane albo odrzucone.

\section{Literatura}

W tej sekcji należy wskazać wykorzystane źródła: książki, artykuły naukowe, dokumentację bibliotek, materiały wykładowe oraz inne źródła użyte podczas pracy nad projektem.

\begin{thebibliography}{9}
\bibitem{sutton_barto}
R. S. Sutton, A. G. Barto, \textit{Reinforcement Learning: An Introduction}, 2nd ed., MIT Press, 2018.
\end{thebibliography}

\end{document}
